\documentclass[10.5pt]{article}
\usepackage{amsmath, amsfonts, amssymb,amsthm}
\usepackage[includeheadfoot]{geometry} % For page dimensions
\usepackage{fancyhdr}
\usepackage{enumerate} % For custom lists
\usepackage{tikz-cd}
\usepackage{graphicx}

\fancyhf{}
\lhead{Lie Algebras Notes}
\rhead{Tighe McAsey}
\pagestyle{fancy}

% Page dimensions
\geometry{a4paper, margin=1in}

\theoremstyle{definition}
\newtheorem{pb}{}
\newtheorem{definition}{Definition}
\newtheorem{proposition}{Proposition}
\newtheorem{example}{Example}
\newtheorem{theorem}{Theorem}
\newtheorem{remark}{Remark}
\numberwithin{proposition}{subsection}
\numberwithin{pb}{subsection}
\usepackage{tikz-cd, stackengine}

% Commands:

\newcommand{\set}[1]{\left\{#1\right\}}
\newcommand{\gen}[1]{\langle#1\rangle}
\newcommand{\abs}[1]{\left\vert#1\right\vert}
\newcommand{\norm}[1]{\lvert\lvert#1\rvert\rvert}
\newcommand{\tand}{\text{ and }}
\newcommand{\tor}{\text{ or }}
\newcommand{\pd}{\frac{\partial}{\partial x_j}}
\newcommand{\px}{\frac{\partial}{\partial x}}
\newcommand{\py}{\frac{\partial}{\partial y}}
\newcommand{\pz}{\frac{\partial}{\partial z}}
\newcommand{\ppx}{\frac{\partial^2}{\partial x^2}}
\newcommand{\ppy}{\frac{\partial^2}{\partial y^2}}
\newcommand{\ppz}{\frac{\partial^2}{\partial z^2}}
\newcommand{\hess}{\operatorname{Hess}}
\newcommand{\inv}{\operatorname{inv}}
\newcommand{\spl}{\operatorname{\mathfrak{sl}}}
\newcommand{\spg}{\mathfrak{g}}
\newcommand{\sph}{\mathfrak{h}}
\newcommand{\spu}{\operatorname{\mathfrak{su}}}
\newcommand{\spo}{\operatorname{\mathfrak{so}}}

\begin{document}
    \section{Heegaard Diagrams of 3-Manifolds}

    \begin{example}[\(\mathbf{L(p,q) := S^3/[(w,z) \sim (\xi_p w, \xi_p^q z)]}\)]

        Consider the standard Heegaard decomposition of \(S^3\) given by 
        \begin{align}
            S^3 = \underbrace{\set{\abs{z} \geq \abs{w}}}_{h^1} \cup_{\underbrace{\abs{z} = \abs{w}}_{\mathbb{T}^2}} \underbrace{\set{\abs{w} \geq \abs{z}}}_{h^2}
        \end{align}
        Then the action of \(\mathbb{Z}/(p)\) which I will denote by \(f\) preserves the handlebodies in this decomposition alongside the torus \(\abs{z} = \abs{w}\). This gives us
        \begin{align}
            L(p,q) = h^1/f \cup_{\mathbb{T}^2/f} h^2/f \label{quot}
        \end{align}
        This is useless unless we can describe what these sets look like, notice that if we can describe \(\mathbb{T}^2/f\), then we have also described the handlebodies, since these can be expressed as \(\mathbb{T}^2/f \times I\). Computing for \(\mathbb{T}^2\):
        \begin{align}
            \mathbb{T}^2/f = \frac{\mathbb{T}^2}{(x,y) \sim \left(x + \frac{1}{p},y+\frac{q}{p}\right)} = \frac{\mathbb{R}^2/(\mathbb{Z}^2)}{(\frac{1}{p},\frac{q}{p})\mathbb{Z}} = \frac{\mathbb{R}^2}{(0,1)\mathbb{Z} \oplus (\frac{1}{p},\frac{q}{p})\mathbb{Z}} \cong \mathbb{T}^2 \label{torusbasis}
        \end{align}
        Thus equation \eqref{quot} also specifies a Heegaard Decomposition. Using the basis computed in equation \eqref{torusbasis} we take the corresponding generators for \(H_1(\partial h^j)\). Now \((0,1)\mathbb{Z}\) is the meridian for \(h^1\) and thus also \(h^1/f\), same for \((1,0)\) as the meridian of \(h^2, h^2/f\), The meridian of \(h^2\) is thus glued to \(p(1/p,q/p) - q(0,1)\). This gives the heegaard decomposition
        on \(\mathbb{T}^2\) as the curves given by \((\alpha, p\beta - q\alpha)\), where \(\alpha,\beta \in H^1(\mathbb{T}^2)\) are the homology basis of the quotient computed in equation \eqref{torusbasis}.
    \end{example}

    \begin{example}[\(\mathbb{T}^3\) with the flat metric]
        Model \(\mathbb{T}^3\) as \([0,2\pi)^3\), then consider
        \begin{align}
            f(x) = \sum_1^3 \frac12 \cos(x^j) + \frac32
        \end{align}
        I claim that this is a self-indexing morse function, checking this
        \begin{align}
            -\nabla f = \frac12 \sin(x^j)dx_j \implies (-\nabla f)^{-1}(0) = \set{(s,t,r) \mid s,t,r \in \set{0,\pi}}
        \end{align}
        And at such points we have
        \begin{align}
            \operatorname{Hess}(f) = \frac12 \operatorname{Diag}(-\cos(x^j))
        \end{align}
        This verifies \(f\) is morse, to see its self indexing note that the index is the number of occurences of \(0\) in the critical point vector.

        Now we solve for a gradient flow of \(f\), from the above computations
        \begin{align}
            \dot{x}^j = \frac12 \sin(x^j)
        \end{align}
        We can apply seperation of variables to \(y'/\sin y = \frac12\) this gives us \(\int \frac{1}{\sin(x^j)} dx^j = \frac{t}{2} + C\) integrating gives us
        \begin{align}
            \log(\tan(x^j/2)) = \frac{t}{2} + C
        \end{align}
        and thus
        \begin{align}
            x^j = 2\arctan(C e^{t/2})
        \end{align}
        Using the initial condition \(x^j(0), C = \tan(x^j(0)/2)\) so that
        \begin{align}
            x^j = 2\arctan(\tan(x^j(0)/2) e^{t/2})
        \end{align}
        The Heegaard diagram uses the asymptotic information of gradient flows, luckily this is easy to compute:
        \begin{align}
            &\lim_{t\to\infty} x^j = \begin{cases}
                \pi & x^j(0) \neq 0 \\
                0 & x^j(0) = 0
            \end{cases} &\lim_{t\to-\infty} x^j = \begin{cases}
                0 &x^j(0) \neq \pi \\
                \pi &x^j(0) = \pi
            \end{cases}
        \end{align}
        (We make sense of this by thinking of \(S^1\) as the codomain of \(\tan\) and domain of \(\arctan\))
        It follows that the stable manifold of the index one critical point \((0,\pi,\pi)\) is given by \(\set{(0,x^2,x^3) \mid x^2,x^3 \neq 0}\). The others are described symmetrically. Similarly we can consider the unstable manifold of the critical point \((\pi,0,0)\) as \(\set{(\pi,x^2,x^3) \mid x^2,x^3 \neq \pi}\), the rest of the cases are symmetric.

        Now, to get the Heegaard decomposition, we study the level set \(f^{-1}(3/2)\), the \(\alpha\) curves are given by
        \begin{align}
            &\alpha_1 = \set{(0,x^2,x^3)\mid \cos(x^2) + \cos(x^3) = -1}\\
            &\alpha_2 = \set{(x^1,0,x^3)\mid \cos(x^1) + \cos(x^3) = -1}\\
            &\alpha_3 = \set{(x^1,x^2,0)\mid \cos(x^1) + \cos(x^2) = -1}
        \end{align}
        And the \(\beta\) curves
        \begin{align}
            &\beta_1 = \set{(\pi,x^2,x^3)\mid \cos(x^2) + \cos(x^3) = 1}\\
            &\beta_2 = \set{(x^1,\pi,x^3)\mid \cos(x^1) + \cos(x^3) = 1}\\
            &\beta_3 = \set{(x^1,x^2,\pi)\mid \cos(x^1) + \cos(x^2) = 1}
        \end{align}
    \end{example}
\end{document}